\subsection{MEAGRE SETS}

DEFINITION 2. A subset A of a topological space X is said to be meagre if it is the union of a countable family of nowhere dense sets.

Equivalently, A is meagre if it is contained in a countable union of closed sets each of which has no interior points.

A meagre set can perfectly well be dense in X; even the whole space X can itself be a meagre set.

An example of the latter possibility is provided by any countable Hausdorff space with no isolated points, e.g., the rational line Q. But a topological space X which is a meagreset in X need not be countable (see Exercise 9).

Every subset of a meagre set in a space X is meagre, and the union of a countable family of meagre sets is meagre.

Let Y be a subspace of X. A subset A of Y is said to be meagre relative to Y if A is meagre when considered as a subset of the topological space Y. It follows from Proposition 2 of no. 1 that if A is a subset of Y which is meagre relative to Y, then A is meagre relative to X; and that if also Y is open in X, every subset A of Y which is meagre relative to X is meagre relative to Y.

